\section{Worker A: target-directed observable-space closure}
\label{sec:linear}

\subsection{The exact fragment}

Let $A_a\in\Q^{D\times D}$ for each $a\in\Sigma$, let $s_0\in\Q^D$, and let
$q\in\Q^{1\times D}$. The problem is
\begin{equation}\label{eq:linear-goal}
 \forall w\in\Sigma^*,\qquad q A_w s_0=0,
 \qquad A_{a_1\cdots a_t}=A_{a_t}\cdots A_{a_1}.
\end{equation}
All words are permitted. A constrained language of inputs needs either a product
with an explicitly supplied finite automaton or a separate proof that the input
restriction is respected. It must not be silently assumed.

Affine updates are included by homogenization:
\[
 s'=M_as+b_a
 \quad\leadsto\quad
 \begin{bmatrix}s'\\1\end{bmatrix}
 =\begin{bmatrix}M_a&b_a\\0&1\end{bmatrix}
  \begin{bmatrix}s\\1\end{bmatrix}.
\]
An affine target is homogenized in the same way. Rational constants remain
rational; no floating-point rank test is needed or appropriate.

For equivalence of two implementations, form their product state and block
updates:
\[
 A_a=\begin{bmatrix}A_a^{(1)}&0\\0&A_a^{(2)}\end{bmatrix},\qquad
 q=\begin{bmatrix}o_1&-o_2\end{bmatrix}.
\]
Initial outputs and every subsequent output are then compared by a single
zero-observation goal. Equal output does not require equal internal states,
equal state dimensions, or the same coordinate system.

\subsection{The algorithm}

Start from the target row. Keep a basis $G$ of the space generated by the rows
seen so far. Whenever an independent row $r$ is admitted, enqueue $rA_a$ for
each action. Dependent rows need not be expanded because multiplying their known
linear combination by $A_a$ produces the corresponding combination of already
scheduled successor rows.

Each queued row carries a concrete forward word. A row representing
$qA_w$ becomes $qA_wA_a$ after right multiplication. Under our execution
convention that represents the word $aw$, not $wa$.

\begin{lstlisting}
queue := [(q, empty_word)]
basis := []
while queue is not empty:
    (r, w) := pop_front(queue)
    if r * initial != 0:
        replay w in the original matrices
        return Counterexample(w)
    if r lies in span(basis):
        continue
    append r to basis
    for action a:
        push_back(queue, (r * A[a], prepend(a, w)))
return exact closure coefficients for the completed basis
\end{lstlisting}

The test of $rs_0$ occurs before discarding a dependent row. It is harmless to
instead perform it after a sound dependence test and propagate the corresponding
zero-value argument; the explicit test makes the counterexample route easier to
audit. The search implementation uses SymPy rational matrices. The checker uses
separately written loops over Python \code{Fraction}, and imports none of the
search code.

\subsection{A certificate without rank computation}

The search returns an $r\times D$ matrix $G$, a row $t\in\Q^{1\times r}$, and
matrices $C_a\in\Q^{r\times r}$ satisfying
\begin{align}
 q&=tG,\label{eq:lin-target}\\
 GA_a&=C_aG && (a\in\Sigma),\label{eq:lin-closure}\\
 Gs_0&=0.\label{eq:lin-init}
\end{align}
The checker verifies these finite identities. It does not check that the rows
of $G$ are independent, that they were obtained by a particular search, or that
$r$ is minimal. Those conditions explain the algorithm's efficiency, but they
are unnecessary for soundness.

\begin{theorem}[Observable certificate soundness]\label{thm:linear-sound}
If \eqref{eq:lin-target}--\eqref{eq:lin-init} hold, then
\eqref{eq:linear-goal} holds.
\end{theorem}
\begin{proof}
The invariant is $Gs=0$. If it holds before action $a$, then
$G(A_as)=(GA_a)s=(C_aG)s=C_a(Gs)=0$. It holds initially by
\eqref{eq:lin-init}, hence after every finite word by induction. Finally,
$qs=tGs=0$ by \eqref{eq:lin-target}.
\end{proof}

This is an inductive invariant, not a numerical approximation to a reachable
set. The certificate can contain redundant rows and still be valid. Conversely,
a claimed invariant basis without \eqref{eq:lin-target} might prove something
irrelevant: the target-membership equation is indispensable.

\subsection{Termination, completeness, and a sharp witness bound}

\begin{theorem}[Finite-dimensional coverage]\label{thm:linear-complete}
Without operational cutoffs, observable closure terminates and decides
\eqref{eq:linear-goal}. If the goal is false and $q\ne0$, a distinguishing word
of length at most $D-1$ exists.
\end{theorem}
\begin{proof}
At most $D$ independent rows can be admitted. Each admitted row creates
$|\Sigma|$ successors, so the queue has at most $1+|\Sigma|D$ entries over the
entire run. On normal termination every basis row is closed under right
multiplication by every action; induction therefore places every $qA_w$ in the
final row space. If all its basis rows annihilate $s_0$, all observations do;
otherwise a tested row supplies a concrete word.

For the length bound let $V_i$ be the span of all target rows arising from words
of length at most $i$. If $V_i=V_{i+1}$, the space is closed under each action
and the chain has stabilized. Starting with the one-dimensional space generated
by nonzero $q$, at most $D-1$ strict dimension increases are possible. Thus
$V_{D-1}$ contains every target row. If its evaluation on $s_0$ is not identically
zero, at least one of its generating words of length at most $D-1$ has nonzero
evaluation.
\end{proof}

The supplied delayed-error fixtures attain this bound. Take a shift matrix with
ones just above the diagonal, initial state $e_D$, and target $e_1^T$. Every
output before step $D-1$ is zero; the next is one. Testing only the first several
inputs would mistake these systems for correct ones. The finite-dimensional
argument, rather than the number of sampled inputs, tells us when a conclusion
is justified.

With a maintained echelon basis, row multiplication and membership reduction
cost $O(D^2)$ field operations per candidate; the search therefore admits an
$O(|\Sigma|D^3)$ implementation. Bit complexity additionally depends on exact
coefficient growth. The prototype recomputes ranks using SymPy rather than
maintaining the optimized incremental representation, so the cubic bound is a
production design target, not a claimed bound for every internal operation of
the delivered Python implementation. This distinction agrees with the
algorithmic refinements in the weighted-automata literature~\cite{kiefer}.

\subsection{A small example that conservation misses}

For input $a\in\{0,1\}$, update
\[
 x'=2x+a,\qquad y'=2y+2a,
 \qquad (x_0,y_0)=(0,0).
\]
The desired observation is $q=2x-y$. Its value changes by the rule $q'=2q$,
not $q'=q$. It is nevertheless zero after every input word. Homogenizing gives
\[
 A_a=\begin{bmatrix}2&0&a\\0&2&2a\\0&0&1\end{bmatrix},\qquad
 G=\begin{bmatrix}2&-1&0\end{bmatrix},\qquad C_a=[2].
\]
The certificate has one invariant row, regardless of word length. Neither
coefficient interpolation over the iteration count nor a guessed closed form
for each state coordinate is necessary.

The corpus also checks product systems related by an invertible triangular
change of coordinates. They are useful because a structural comparison of the
two state vectors fails, whereas observable closure automatically discovers
precisely the linear combinations needed to establish output equivalence.

\subsection{Beyond a finite alphabet: a careful extension}
\label{sec:symbolic-input}

Suppose an update depends polynomially on a fresh scalar input $u$:
\[
 A(u)=\sum_{j=0}^d u^j M_j.
\]
A sufficient universal certificate consists of
\[
 GM_j=C_jG\qquad(0\le j\le d),
\]
because then $GA(u)=(\sum_j u^jC_j)G$ for every $u$. Over an infinite field,
vanishing of an output polynomial for all choices of inputs forces all its
coefficients to vanish, so coefficient closure gives a completeness route for
this precise polynomial-input linear-state fragment.

A coefficient matrix is not an input symbol, however. If a coefficient row does
not annihilate the initial state, its index sequence is only evidence of a
nonzero output polynomial. To produce an actual counterexample, instantiate
that polynomial on an exact grid whose coordinate sizes exceed the respective
degrees, then replay the chosen rational inputs. Presenting the coefficient
indices as user inputs would be a semantic bug. This extension is designed here
but not implemented in the recorded corpus.

Finite-alphabet closure is the recommended first Lean milestone. Polynomial
input lifting should follow only after the source bridge, coefficient identity,
and concrete-input reconstruction have their own tests.
